\documentclass[11pt,a4paper]{article}
\usepackage[UTF8]{ctex}
\usepackage{amsmath, amssymb, amsthm, amsfonts}
\usepackage{geometry}
\usepackage{tcolorbox}
\usepackage{tasks}

\geometry{left=2cm, right=2cm, top=2.5cm, bottom=2.5cm}

\newtcolorbox{mybox}[1]{
    colback=white,
    colframe=black,
    coltitle=black,
    fonttitle=\bfseries\Large\linespread{1.5}\selectfont,
    sharp corners,
    boxrule=0.8pt,
    toptitle=3mm,
    bottomtitle=3mm,
    top=4mm,
    bottom=4mm,
    left=2mm,
    right=2mm,
    titlerule=0.5pt,
    attach title to upper,
    after title={\par\medskip\hrule\medskip},
    title=#1
}

\settasks{
    label=\Alph*.,
    label-width=1.5em,
    item-indent=2em,
    column-sep=10pt
}

\begin{document}

\begin{mybox}{单选题：集合的交集运算}

已知集合 \( A = \{ x \mid 1 < x^2 < 9 \} \)，\( B = \{-2, -1, 0, 1, 2\} \)，则 \( A \cap B = ( ) \)

\begin{tasks}(4)
    \task \( \{-2, -1, 0\} \)
    \task \( \{-2, 2\} \)
    \task \( \{-1, 0, 1\} \)
    \task \( \{0, 1, 2\} \)
\end{tasks}

\end{mybox}

\vspace{2em}

\begin{mybox}{填空题：函数定义域}

函数 \( f(x) = \sqrt{4 - x^2} + \ln(x + 1) \) 的定义域为 \underline{\hspace{4cm}}.

\end{mybox}

\end{document}
